\section{Background: the failure we had to handle}
\label{sec:background}

\subsection{How the RPC backend distributes a model}

Built with \code{-DGGML\_RPC=ON}, \texttt{llama.cpp} provides
\code{ggml-rpc-server}, which advertises the accelerators of the machine it runs
on. A client passes \code{--rpc host:port} and those remote devices join the
local ones as offload targets. Model layers are then distributed across the
combined set, and \code{--tensor-split} can override the default proportions.

Two properties of this mechanism shaped our design.

First, the device order used by \code{--tensor-split} is not the order printed
by \code{--list-devices}. We established the real order by measurement rather
than by reading documentation, loading the same 3B model twice with opposite
weights:

\begin{center}\small
\begin{tabular}{@{}lrr@{}}
\toprule
\code{--tensor-split} & ROG CUDA0 & TUF RPC0 \\
\midrule
\code{1,4} & \SI{1573}{MiB} & \SI{345}{MiB} \\
\code{4,1} & \SI{552}{MiB}  & \SI{1366}{MiB} \\
\bottomrule
\end{tabular}
\end{center}

Weighting the first position moved memory onto the local GPU, so RPC devices are
listed first and the local device last. Our initial implementation had this
backwards, which placed most of the model on the controller while the logs
claimed it was distributed. Reading the placement back from
\texttt{llama.cpp}, rather than trusting the request, is what exposed the error.

Second, the RPC server accepts one client at a time. While a model is resident,
the device cannot be queried by a second process. We confirmed this when a
device probe aborted against an endpoint already held by the running model. This
rules out querying remote devices for live memory figures during serving, and it
is why worker-reported telemetry is the only usable source.

\subsection{What happens when a worker disappears}

The behaviour on peer loss is unambiguous. With a 3B model split across the
Linux and Windows laptops, we began a request and then stopped the remote
\code{ggml-rpc-server} by hand. \texttt{llama.cpp} emitted:

\begin{quote}\small\ttfamily
Remote RPC server crashed or returned malformed response
\end{quote}

\noindent
and the process aborted. Measured from the client side: \code{curl} exited 52
with HTTP status 000 after \SI{34.06}{\second}, no response body was produced,
the process was gone five seconds later, the health endpoint was unreachable,
and nothing restarted. The model had to be loaded again from scratch.

This is reasonable behaviour for a proof-of-concept backend. A lost peer means
part of the computation graph is unreachable mid-forward-pass, and aborting is
safer than continuing with an incomplete model. The consequence, though, is that
the reliability of the whole pool equals the reliability of its least stable
member. For laptops on Wi-Fi, that is low.

\subsection{The resulting requirement}

Recovery cannot be implemented inside the forward pass without modifying
\texttt{llama.cpp} itself, which we chose not to do: it would fork a fast-moving
upstream project and make the result harder to reproduce. The alternative is to
treat the model process as replaceable and to supervise it from outside. That
places three obligations on the supervisor:

\begin{itemize}[leftmargin=1.4em,itemsep=2pt]
  \item \textbf{Distinguish failure modes.} An ordinary bad request must not
        trigger a rebuild, and a dead worker must not be mistaken for one. We
        settled on process exit as the signal, for reasons given in
        Section~\ref{sec:design}.
  \item \textbf{Verify endpoints independently of agent health.} A worker's
        heartbeat proves its agent is running, not that its accelerator is
        reachable. The experiment in Section~\ref{sec:recovery} was designed
        specifically to separate these two states.
  \item \textbf{Preserve the interrupted request.} A user who submitted a
        prompt before the failure should not have to resubmit it.
\end{itemize}
